% !TEX root = chapter-standalone.tex

\section{Design and co-design}
\label{sec:intro-design}

In this chapter we put the order theory of the previous chapters to work.
We will see that \SY{preorders}, \SY{monotone functions}, and monotone relations provide a simple and surprisingly general language for describing design, and in particular for describing how the design of a system depends on the design of its components.
Before we introduce any mathematics, this section explains, informally, what we mean by design and co-design.

\subsection{What is design?}

We take a broad view of what it means to design, one that is not limited to engineering.
In the words of Herbert Simon (\cite{hebert96sciences}, Chapter 5):

\begin{quote}
    Everyone designs who devises courses of action aimed at changing existing situations into preferred ones.
    The intellectual activity that produces material artifacts is no different fundamentally from the one that prescribes remedies for a sick patient or the one that devises a new sales plan for a company or a social welfare policy for a state.
\end{quote}

The examples in this book are mostly taken from engineering, because it is easy to imagine building a machine out of simpler components, and to imagine the choices and trade-offs involved.
But the ideas apply equally well to other disciplines, as long as we are willing to take an abstract view of what a ``system'' and a ``component'' are.
In urban planning, the components of a city are roads, sewers, and residential areas; in economics, a component might be an incentive scheme designed to move a market towards a more desirable state.

\subsection{Functionalities and requirements}

Imagine that you are an engineer standing in front of an empty whiteboard, about to design a new product.
The first question to ask is: what is the product \emph{for}?
The answer is expressed by its \emph{functionalities}: what the product must provide, deliver, or achieve.
For example:
\begin{itemize}
    \item a car must be able to transport at least four passengers;
    \item a battery must store at least $\unit[100]{kJ}$ of energy;
    \item a delivery drone must be able to carry a parcel of at least $\unit[2]{kg}$.
\end{itemize}
(Engineers often say ``function'' here, but since that word already has a precise mathematical meaning, we will always say ``functionality''.)

The second question is: what does the product \emph{need}?
The answer is expressed by its \emph{requirements}: the resources that must be spent or made available in order to provide the functionalities.
For example:
\begin{itemize}
    \item the car should not cost more than $\unit[30000]{CHF}$;
    \item the battery should not weigh more than $\unit[1]{kg}$;
    \item the drone should not need more than $\unit[10]{min}$ to charge.
\end{itemize}

There is a pleasant duality between the two.
Functionalities are typically given as \emph{lower bounds} (``at least''), and requirements as \emph{upper bounds} (``at most'').
Most design objectives can be phrased as either providing as much functionality as possible, or needing as few resources as possible.

Two features of functionalities and requirements will be central in what follows.
First, they come with a natural notion of ``less'' and ``more'', and this notion is often only partial.
A battery that stores more energy but also weighs more is neither better nor worse than a lighter battery that stores less; the two are incomparable, and choosing between them is a matter of trade-offs.
This is why we model functionalities and requirements by \SY{preorders} and \SY{posets} (\cref{chap:intro-tradeoffs}).
Second, the relationship between functionalities and requirements is \emph{monotone}.
If some requirements suffice to provide a functionality, then they also suffice to provide any lesser functionality; and if a functionality can be provided with some requirements, then it can also be provided with larger ones.
To obtain more, you need more.

\subsection{Everything is a resource}

Here is a further observation that shapes the whole theory.
Consider a small electric delivery drone.
Its motors provide thrust, and they require electric power.
The electric power is provided by a battery, which in turn requires a certain mass (and money, and charging time).
The thrust, in turn, is required by the drone as a whole in order to lift the parcel.

In this chain, the electric power is a \emph{requirement} of the motors, but it is a \emph{functionality} of the battery.
Similarly, thrust is a functionality of the motors, but a requirement of the drone.
So ``functionality'' and ``requirement'' are not two different kinds of things: they are two \emph{roles} that a resource can play for a given component.
Every poset that appears in the theory is a poset of resources, and each component relates the resources it provides, its functionalities, to the resources it needs, its requirements.
Components are connected precisely by letting the requirements of one component be provided as functionalities by another.
This is why we will always speak of a component as relating \emph{functionalities} to \emph{requirements}, while keeping in mind that both are resources.

\subsection{From components to systems: co-design}

A system is composed of components.
The computer pioneer Howard Aiken is credited with the remark that a system is composed of components, and a component is something you understand.
For design, we modify this slightly:

\begin{quote}
    A system is composed of components;\\
    a component is something you understand \emph{how to design}.
\end{quote}

The difficulty is that the components of a system cannot be designed one at a time.
In the drone, a larger battery provides more power, but it is also heavier, so the drone needs more thrust, so the motors need more power, which calls for a larger battery.
The requirements of one component depend on choices made for other components, sometimes in a circular way.
Choices made component by component, without looking at the system as a whole, are therefore likely to be suboptimal, or even infeasible.

We use the word \emph{co-design} for the activity of making the design choices for all components of a system \emph{together}, so as to meet the system-level functionalities with the smallest possible requirements.
The prefix ``co'' has several meanings for us.
It stands for \emph{compositional}: the design of the system is obtained by composing the designs of its components.
It stands for \emph{collaborative}: in practice, different components are designed by different people or teams, who need a common language to describe what each component provides and requires.
(Melvin Conway famously observed that organizations tend to produce systems whose structure mirrors their own communication structure, an observation that has since been studied empirically \cite{maccormack12exploring}.)
And it stands for \emph{computational}: we want descriptions of components that are precise enough for a computer to find optimal designs.

\subsection{What we model in this chapter}

In this chapter we take the first steps towards a mathematical theory of co-design.
\begin{itemize}
    \item We model a single component by a \emph{feasibility relation}, or equivalently a \emph{design problem}, which says, for each functionality and each requirement, whether the functionality can be provided with that requirement (\cref{sec:dpdefinition}).
          The monotonicity described above is built into the definition.
          We do not model \emph{how} a functionality is provided, only \emph{whether} it can be.
    \item We connect components in series, letting the requirement of one component be provided as a functionality by another (\cref{sec:dp-series-composition}).
          Other ways of connecting components, including the circular dependencies of the drone example, will be studied later in the book.
    \item We ask the questions that designers ask of a component or a system: given a functionality, which requirements suffice, and which are minimal? Given a requirement, which functionalities can be provided (\cref{sec:dp-querying})?
    \item Finally, we see that the answers to such questions are, in general, not single optimal choices but collections of incomparable optimal trade-offs, described mathematically by \SY{antichains} (\cref{sec:dp-optimality}).
\end{itemize}
